\documentclass[a4paper,12pt]{report}
\usepackage[utf8]{inputenc}
\usepackage[T1]{fontenc}
\usepackage[french]{babel}
\usepackage{graphicx}
\usepackage{hyperref}
\usepackage{geometry}
\usepackage{titlesec}
\usepackage{fancyhdr}
\usepackage{float}
\usepackage{listings}
\usepackage{color}

% Configuration des marges
\geometry{hmargin=2.5cm,vmargin=2.5cm}

% Configuration de l'entête et pied de page
\pagestyle{fancy}
\fancyhf{}
\rhead{Traitement d'Images et SIG}
\lhead{GeoProd\_CM}
\cfoot{\thepage}

% Informations du document
\title{\textbf{Développement d'une Plateforme SIG Web pour la Visualisation des Bassins de Production au Cameroun}}
\author{Étudiants du Master Informatique}
\date{\today}

\begin{document}

% Page de garde
\begin{titlepage}
    \centering
    \includegraphics[width=0.3\textwidth]{backend/static/img/logo.png}\\[1cm]
    {\scshape\LARGE Université de Yaoundé 1 \par}
    {\scshape\Large Faculté des Sciences \par}
    {\scshape\Large Département d'Informatique \par}
    \vspace{1cm}
    {\Large Cours de Traitement d'Images et SIG \par}
    \vspace{0.5cm}
    {\large Enseignant : Dr Hyppolyte TAPAMO \par}
    \vspace{2cm}
    \rule{\linewidth}{0.5mm} \\[0.4cm]
    {\huge\bfseries Plateforme Web SIG : GeoProd\_CM \par}
    {\Large Visualisation et Analyse des Données de Production Agro-Pastorale \par}
    \rule{\linewidth}{0.5mm} \\[1.5cm]
    
    {\large\itshape Rapport de Projet Académique \par}
    \vspace{2cm}
    \vfill
    {\large Année Académique 2025-2026 \par}
\end{titlepage}

\tableofcontents
\newpage

\chapter{Introduction}
\section{Contexte et Problématique}
Le Cameroun, souvent qualifié de "Grenier de l'Afrique Centrale", repose économiquement sur son secteur primaire (agriculture, élevage, pêche). Cependant, les données statistiques relatives à ces productions restent souvent dispersées dans des rapports statiques (PDF, Excel), rendant difficile leur analyse spatiale et décisionnelle.

Les décideurs et acteurs économiques manquent d'outils interactifs pour visualiser rapidement :
\begin{itemize}
    \item Où se situent les bassins de production majeurs ?
    \item Comment évolue la production par région ou département ?
    \item Quelles sont les zones dominantes pour un produit spécifique (Cacao, Café, Maïs...) ?
\end{itemize}

\section{Objectifs du Projet}
L'objectif de ce projet est de concevoir et déployer une plateforme Web SIG (\textit{Système d'Information Géographique}) capable de :
\begin{enumerate}
    \item Centraliser les données géographiques (découpage administratif) et statistiques.
    \item Offrir une visualisation cartographique interactive (Choroplèthe).
    \item Permettre l'exportation et le filtrage dynamique des données.
\end{enumerate}

Le projet est accessible publiquement à l'adresse : \url{https://geoprod-cm.onrender.com}.

\chapter{Méthodologie et Analyse}
\section{Modélisation des Données}
Nous avons structuré notre base de données relationnelle autour de quatre entités principales pour refléter la hiérarchie administrative du Cameroun :
\begin{itemize}
    \item \textbf{Region} : Les 10 régions du pays.
    \item \textbf{Departement} : Subdivisions des régions.
    \item \textbf{Arrondissement} : Unités administratives locales.
    \item \textbf{Production} : Table factuelle liant une quantité, un produit, une année et une zone géographique.
\end{itemize}

\section{Architecture Technique}
Le projet suit une architecture \textbf{Client-Serveur} moderne :
\begin{itemize}
    \item \textbf{Backend (API)} : Développé avec \textbf{Django} et \textbf{Django REST Framework}. Il expose les données au format JSON pour l'application frontend.
    \item \textbf{Base de Données} : \textbf{PostgreSQL} a été choisi pour sa robustesse et sa capacité à gérer des données spatiales (via PostGIS ou JSONFields).
    \item \textbf{Frontend (UI)} : Une interface responsive construite avec HTML5, \textbf{Tailwind CSS} pour le style, et \textbf{Leaflet.js} pour le moteur cartographique.
\end{itemize}

\chapter{Implémentation Technique}

\section{Visualisation Spatiale avec Leaflet}
L'un des défis majeurs a été la représentation des données sur la carte. Nous avons opté pour une carte \textbf{Choroplèthe} (carte par plages de couleurs).
L'algorithme implémenté côté client :
\begin{enumerate}
    \item Récupère les productions filtrées par l'utilisateur.
    \item Calcule dynamiquement les seuils (min/max) de production.
    \item Attribue une intensité de couleur (du vert clair au vert foncé) à chaque zone géographique en fonction de sa production relative.
\end{enumerate}

\section{Optimisation des Performances}
Pour garantir une expérience fluide malgré le volume de données :
\begin{itemize}
    \item \textbf{Backend} : Utilisation de \texttt{select\_related} dans les requêtes Django pour éviter le problème de "N+1 query".
    \item \textbf{Pagination} : Mise en place d'une pagination serveur (20 items/page) pour le tableau de données.
    \item \textbf{Fichiers Statiques} : Utilisation de \textbf{Whitenoise} avec compression Gzip/Brotli pour servir les assets rapidement.
\end{itemize}

\section{Déploiement avec Docker}
Afin d'assurer la portabilité et la reproductibilité du projet, nous avons "conteneurisé" l'application avec **Docker**.
\begin{itemize}
    \item Image de base : \texttt{python:3.11-slim} (légère et sécurisée).
    \item Script d'entrée (\texttt{entrypoint.sh}) : Automatise les migrations de base de données au démarrage du conteneur.
    \item Serveur WSGI : \textbf{Gunicorn} est utilisé pour servir l'application Python en production.
\end{itemize}

\chapter{Résultats et Fonctionnalités}

\section{Interface Cartographique}
L'utilisateur peut filtrer par Secteur (Agriculture, Élevage), Produit et Année. La carte se met à jour instantanément.
\begin{figure}[H]
    \centering
    % DECOMMENTER ET AJOUTER L'IMAGE DANS LE DOSSIER
    % \includegraphics[width=1\textwidth]{backend/static/img/screenshot_map.png}
    \caption{Interface principale : Carte interactive des productions de Cacao dans le Centre.}
\end{figure}

\section{Tableau de Bord et Analyse}
Une vue détaillée permet de consulter les données brutes, avec une synthèse statistique (Total, Zone dominante) et un module d'export Excel.
\begin{figure}[H]
    \centering
    % DECOMMENTER ET AJOUTER L'IMAGE DANS LE DOSSIER
    % \includegraphics[width=1\textwidth]{backend/static/img/screenshot_data.png}
    \caption{Page de données : Synthèse statistique et tableau filtrable.}
\end{figure}

\chapter{Conclusion et Perspectives}
Ce projet académique a permis de réaliser une chaîne complète de traitement de l'information géographique, de la base de données à la visualisation web. La plateforme \textit{GeoProd\_CM} est fonctionnelle, performante et déployée.

\section*{Perspectives d'amélioration}
\begin{itemize}
    \item Intégration de données météorologiques en temps réel.
    \item Développement d'une application mobile pour la collecte de données sur le terrain.
    \item Ajout d'analyses prédictives (Machine Learning) sur les rendements agricoles.
\end{itemize}

\section*{Remerciements}
Nous adressons nos sincères remerciements au \textbf{Dr Hyppolyte TAPAMO} pour la qualité de son enseignement en Traitement d'Images et SIG, ainsi que pour son accompagnement tout au long de ce projet.

\end{document}
